\chapter{Calibración}
\label{ch:calibracion}

La calibración es \textbf{el primer paso} del análisis. Determina la relación
entre canal y velocidad (mm/s) y, opcionalmente, la referencia de desplazamiento
isomérico. En Fitbauer la calibración se \textbf{fija} y se reutiliza en todas
las cargas posteriores, de modo que no hay que redefinirla para cada espectro.

\section{Por qué hace falta calibrar}

El transductor de velocidad recorre un ciclo triangular; el sistema registra las
cuentas por canal. Para convertir el eje de canales en un eje físico de
velocidad se necesita la \textbf{velocidad máxima} \param{Vmax} (mm/s en el
extremo del recorrido). Esta \param{Vmax} depende de la amplitud del transductor
y es común a todos los espectros medidos con la misma configuración: por eso se
mide una vez con un patrón y se reutiliza.

El patrón habitual es el \textbf{$\alpha$-Fe metálico}, cuyo sexteto tiene
posiciones de línea muy bien conocidas. Fitbauer adopta la convención de que un
espectro de $\alpha$-Fe corresponde a $\BHF = 33{,}0$~T, usando el patrón de
velocidad publicado ($\pm0{,}839 / \pm3{,}084 / \pm5{,}329$~mm/s), igual que
NORMOS.

\begin{aviso}
Las posiciones del sexteto de referencia \textbf{no} se derivan de los momentos
nucleares: el cálculo de libro da un desdoblamiento $\sim0{,}4\%$ menor y haría
que el $\BHF$ saliera $\sim0{,}1$~T demasiado alto. Se usan las posiciones
publicadas de $\alpha$-Fe a 33,0~T.
\end{aviso}

\section{Definir una calibración a partir de un $\alpha$-Fe}
\label{sec:definir-calib}

El procedimiento habitual es cargar un espectro de $\alpha$-Fe, ajustarlo como
sexteto (véase el capítulo~\ref{ch:simular}) y marcarlo como calibración:

\begin{enumerate}[nosep]
  \item \menu{Archivo > Cargar…} y seleccione el espectro de $\alpha$-Fe
        (por ejemplo \ruta{data\_sample/hierro\_metalico\_alphaFe.adt}).
  \item Ajuste el sexteto (\menu{Ajuste > Ajustar}) hasta reproducir las seis
        líneas. Puede activar \emph{Ajustar Vmax} para refinar la propia
        \param{Vmax} contra las posiciones de $\alpha$-Fe.
  \item Marque el resultado como calibración con
        \menu{Archivo > Usar fichero actual como calibración…}, o desde el menú
        contextual del cuadro de fichero (clic derecho).
\end{enumerate}

\autofig{fig_calibration_fit}{Ajuste de un espectro de $\alpha$-Fe usado como
calibración. El $\BHF$ recuperado ($\approx 33$~T) y el $\chi^2$ reducido
próximo a 1 indican una calibración correcta. Figura generada con
\ruta{scripts/make\_figures.py}.}

Al marcar la calibración, Fitbauer ofrece dos variantes desde el menú
contextual del cuadro de fichero:

\begin{itemize}[nosep]
  \item \menu{Usar como calibrado (vmax e iso actuales)}: toma directamente la
        \param{Vmax} y el desplazamiento isomérico $\delta$ del estado actual.
  \item \menu{Usar como calibrado (detalle)…}: abre un diálogo para revisar y
        editar el nombre de la muestra, la \param{Vmax} y el IS antes de fijar.
\end{itemize}

\guicap{gui_use_as_calibration}{Diálogo «Usar como calibración» con el nombre de
la muestra, la Vmax calibrada y el desplazamiento isomérico de referencia.}

\section{La calibración queda \emph{fijada}}
\label{sec:calib-fija}

Una vez definida, la calibración queda \textbf{fijada}: su \param{Vmax} (y el IS
de referencia) se guardan y se muestran en la etiqueta de calibración con la
marca \textbf{«fija»}. A partir de ese punto se aplica automáticamente a los
espectros que cargue, según esta regla:

\begin{center}
\begin{tabular}{@{}p{0.46\textwidth}p{0.46\textwidth}@{}}
\toprule
\textbf{Acción} & \textbf{Efecto sobre la calibración} \\
\midrule
Cargar un fichero \textbf{sin} calibración propia (cuentas en bruto:
\ruta{.ws5}, \ruta{.adt}) &
Se \textbf{usa la fija}: el eje de velocidad se construye con la \param{Vmax}
fijada. La calibración \textbf{no varía}. \\[6pt]
Cargar un fichero \textbf{con} calibración propia (eje de velocidad:
\ruta{.csv}, \ruta{.txt}, \ruta{.dat}, \ruta{.exp}) &
\textbf{Sustituye} la calibración por la del fichero (su \param{Vmax}) y se
\textbf{avisa} en la barra de estado. \\
\bottomrule
\end{tabular}
\end{center}

Es decir: si fija una calibración y después carga \textbf{solo datos} de cuentas
en bruto, la calibración permanece \textbf{intacta} y se aplica a todos ellos.

\begin{nota}
Los ficheros de cuentas en bruto (\ruta{.ws5}/\ruta{.adt}) no llevan un eje de
velocidad; necesitan la \param{Vmax} de la calibración para construirlo. Los
ficheros ya doblados en velocidad (\ruta{.csv}, etc.) traen su propio eje: por
eso \emph{sustituyen} la calibración activa.
\end{nota}

\guicap{gui_calib_label_fixed}{Etiqueta de calibración mostrando el estado
«fija», con la muestra, la Vmax y el IS de referencia.}

\subsection{Sustitución de la calibración}

Cuando carga un fichero con calibración propia estando ya fijada otra, la nueva
sustituye a la anterior y Fitbauer muestra un aviso temporal en la barra de
estado del tipo:

\begin{center}\ttfamily\small
Calibración sustituida por la del fichero datos.csv: Vmax 12.0070 $\to$ 11.5000 mm/s
\end{center}

Así queda constancia de que la \param{Vmax} activa ha cambiado. El
desplazamiento isomérico de referencia se conserva.

\subsection{Liberar la calibración}

Para dejar de aplicar la calibración fijada, use \menu{Liberar calibración
fijada} en el menú contextual del cuadro de fichero. A partir de entonces cada
carga usará la \param{Vmax} que corresponda a su propio fichero (o el valor por
defecto del control).

\section{Calibraciones desde la web}

Si dispone de acceso a la API, \menu{Archivo > Web > Calibraciones web…} permite
descargar una calibración remota. Sus valores (\param{velocity\_calibrated} e
\param{isomer\_shift}) quedan fijados exactamente igual que una calibración
local, reescalando el eje de velocidad del espectro actual.

\section{Persistencia en la sesión}

La calibración fijada forma parte de la \textbf{sesión}: al guardar
(\menu{Archivo > Guardar sesión…}) se almacena junto con el resto del estado, y
al cargarla (\menu{Archivo > Cargar sesión…}) se restaura tal cual. De este modo
un análisis es reproducible sin volver a calibrar.

\begin{truco}
Para un trabajo por lotes sobre muchas muestras medidas con la misma
configuración: cargue y fije la calibración de $\alpha$-Fe una vez, guárdela en
una sesión base y arranque desde ella cada día.
\end{truco}
